% !TeX root = ../main.tex

Dieses Kapitel beschreibt das Konzept der Experimente zur Untersuchung verschiedener Retry Strategien für
Transaktionen in PostgreSQL\@.
Ziel ist es, einen reproduzierbaren Versuchsaufbau zu definieren, mit dem das Verhalten der Retry-Mechanismen in
fest erzeugten Fehlersituationen analysiert werden kann.

Dafür wird zunächst die Zielsetzung erläutert.
Anschließend wird die Systemarchitektur beschrieben und wie sie sich mit PostgreSQL ergänzen.
Darauf aufbauend wird der Aufbau des Retry-Mechanismus vorgestellt, sowie der Ablauf eines vollständigen
Experiments erläutert.

\section{Zielsetzung}
\label{sec:purpose}
% !TeX root = ../../main.tex

Das Ziel der Experimente ist es die einzelnen, Retry systematisch zu testen und zu untersuchen, wie dieser sich unter
den verschiedenen Konfigurationen verhält.
Zusätzlich soll analysiert werden wie stark der Einfluss der verschiedenen Parameter ist.

Die Parameter sind die Anzahl an erlaubten Retries, \emph{Retry Count}, die Dauer des Wartens zwischen den einzelnen
Retries, \emph{Delay}, und die einzelnen Strategien.
Die Strategien, wie in Abschnitt~\ref{sec:retry} erwähnt, sind als Basis kein Retry, immediate Retry, Retry mit festem
Delay und der Retry mit exponential backoff + jitter Delay.

Bewertet werden soll das mit verschiedenen Kenngrößen.
Die wichtigste hierbei ist die Erfolgsrate, die sich berechnet, in dem alle erfolgreichen Transaktionen durch die
Gesamtanzahl an Transaktionen geteilt werden.
Eine weitere entscheidende Kenngröße ist die durchschnittliche Ausführungszeit der Transaktionen.
Die berechnet sich aus der Summe der \texttt{total\_ms} aller Transaktionen geteilt durch die Anzahl der Transaktionen.
Des Weiteren werden die Quartile berechnet.
Dabei wird der Fokus auf die 50\%, 95\%, 99\% und 99.9\% Quartile gelegt, um zu zeigen, wie viel Prozent der
Transaktionen bis zu welchem Zeitpunkt abgeschlossen haben.
Die \emph{retry\_distribution} zeigt, wie viele Transaktionen wie häufig einen Retry ausgeführt haben.
Sortiert wird bei den genannten Kenngrößen nach successful oder failed.

Eine weitere Kenngröße ist der \emph{retry\_overhead}, der sich durch die durchschnittliche Differenz zwischen der
Gesamtdauer \texttt{total\_ms} und dem besten Versuch \texttt{best\_ms} berechnet.
\[
	\mathrm{retry\_overhead} =
	\frac{1}{n}
	\sum_{i=1}^{n}
	\left(
		\mathrm{total}_{\mathrm{ms},i}
		-
		\mathrm{best}_{\mathrm{ms},i}
	\right),
	\quad n = \mathrm{total\_transactions}
\]

Dieser Wert ist besonders relevant, da er den zusätzlichen Zeitaufwand zeigt, der durch die Ausführung von Retries
entsteht.
Dadurch lässt sich bewerten, welchen Einfluss Wiederholungsversuche auf die Gesamtausführungszeit einer Transaktion
haben.


Die letzte Kenngröße ist mit, eine der wichtigsten, der \emph{throughput}.
Der Durchsatz zeigt an, wie viele erfolgreiche Transaktionen innerhalb einer bestimmten Zeit verarbeitet werden
können.
Dadurch ermöglicht er eine Bewertung der Leistungsfähigkeit des Systems.
Dieser berechnet sich aus der Anzahl erfolgreicher Versuche(N) geteilt durch die Dauer des Experiments.
\[
	\mathrm{throughput} =
	\frac{
		N_{\mathrm{success}}
	}{
		t_{\mathrm{finish}}-t_{\mathrm{start}}
	}
\]
Es wird durch die Dauer des Experiments geteilt, da diese fest definiert sind.
In den Experimenten sind die Parameter variable, aber wenn ein Experiment angefangen wurde, läuft das die Anzahl an
Iterationen durch.
Dadurch lässt sich der Durchsatz der einzelnen Experimente berechnen und durch weitere Sortierungsparameter wird der
Durchsatz für diesen Parameter berechnet.

Damit diese Experimente repräsentative Ergebnisse liefern, gibt es Rahmenbedingungen.
Die Fehlerszenarien sind fest definiert für alle Transaktionen.
Wenn Retry beispielsweise für Deadlocks getestet werden soll, dann erzeugen die Transaktionen auf jeden Fall auch
einen Deadlock.
Bei den anderen Fehlerklassen passiert genau das Gleiche, dadurch lässt sich das Verhalten unter reproduzierbaren
Bedingungen beobachten.

Die Kenngrößen sollen im Endeffekt zeigen, ob es Retry Strategien gibt, die bei den einzelnen Fehlerklassen bessere
Ergebnisse liefern als andere Strategien.
Zusätzlich soll anhand der Ergebnisse abgeleitet werden, wann Retry sinnvoll eingesetzt werden kann.

\section{Gesamtarchitektur}
\label{sec:architecture}
% !TeX root = ../../main.tex

Bitti hol mich hier raus!!!

\section{Aufbau}
\label{sec:design}
% !TeX root = ../../main.tex

Die Experimente bestimmen, mit welcher Konfiguration eine Retry Strategie ausgeführt wird.
Eine Konfiguration besteht aus den folgenden Parametern.
\begin{itemize}
	\item Retry Count
	\item Retry Delay
	\item Retry Strategie
	\item Concurrency
	\item Workload (nicht variabel)
	\item Iterations (nicht variabel)
\end{itemize}
Der Retry Count legt fest, wie häufig eine Transaktion wiederholt werden darf.
Retry Delay beschreibt, wie lange eine Strategie initial warten muss, bevor sie die Transaktion wiederholen darf.
Durch Retry Strategie legt ein Experiment fest, mit welchen Strategien es den Retry ausführen soll.
Der Concurrency Parameter legt fest, wie viele gleichzeitige Anfragen eine Datenbank bekommt, die alle den SQL-Code
ausführen wollen.
Workload ist fest definiert, damit man vergleichbare Ergebnisse bekommt und Iterations sorgt dafür, dass die
Ergebnisse eine gewisse Stabilität haben und keine großen Ausreißer sind.
Durch diese Parameter lassen sich die Strategien testen, ohne dass es zu viel wird.

Der Aufbau einer Retry Strategie besteht daraus, dass sie SQL-Code ausführen soll.
Daraufhin bekommt sie von der Datenbank einen Status zurück, entweder ist die Transaktion erfolgreich gewesen oder
nicht.
Im nächsten Schritt wird dies getestet.
Gleichzeitig wird noch getestet ob der Retry Count überschritten wird oder nciht.
Ist die Überprüfung positiv, dass ein weiterer Retry möglich ist, wird der Retry angewandt.
Es wird der Retry Count erhöht und so lange gewartet, wie Delay besagt.
Danach wird die Transaktion erneut ausgeführt.

Dies geschieht so lange bis eine der Abbruchbedingungen zutrifft.
Entweder war die Transaktion erfolgreich, oder die Transaktion hat die maximal zulässige Anzahl an Wiederholungen
überschritten.
Ist das passiert, wird das Ergebnis zurückgegeben und in die Logdatei geschrieben.

\section{Ablauf des Systems}
\label{sec:process}
% !TeX root = ../../main.tex

Abbildung~\ref{fig:experiment_flow} veranschaulicht den grundlegenden Ablauf eines Experiments.
Nach dem Start der Anwendung werden zunächst die für das jeweilige Experiment definierten Parameter geladen
und die entsprechenden Versuchsbedingungen festgelegt.
Dazu gehören unter anderem die Anzahl der durchzuführenden Transaktionen sowie die für den Versuch relevanten
Konfigurationsparameter.
Anschließend wird das Experiment gestartet und die einzelnen Transaktionen entsprechend der vorgegebenen Konfiguration
ausgeführt.
Für jede Transaktion wird zunächst eine SQL-Anfrage an die Datenbank gesendet und ausgeführt.
Nach der Ausführung wird das Ergebnis der Transaktion ausgewertet.
Dabei wird insbesondere geprüft, ob die Transaktion erfolgreich abgeschlossen wurde oder ob ein Fehler aufgetreten ist,
der einen erneuten Ausführungsversuch erforderlich macht.
Tritt ein entsprechender Fehler auf und steht gemäß der definierten Retry-Strategie noch ein weiterer Versuch zur
Verfügung, wird die Transaktion erneut ausgeführt.
Dieser Ablauf kann sich abhängig vom Ergebnis und der maximal zulässigen Anzahl an Wiederholungsversuchen
mehrfach wiederholen.
Sobald die Transaktion erfolgreich abgeschlossen wurde oder keine weiteren Versuche mehr möglich sind, wird mit der
nächsten Transaktion fortgefahren.
Nach Abschluss aller für das Experiment vorgesehenen Transaktionen werden die während der Ausführung erfassten
Ergebnisse und Messwerte gespeichert.
Hierzu gehören hauptsächlich die Informationen über erfolgreiche und fehlgeschlagene Transaktionen sowie die Anzahl
der durchgeführten Wiederholungsversuche.
Auf Grundlage dieser aufgezeichneten Daten werden anschließend die für die Auswertung relevanten Kenngrößen berechnet.
Dadurch kann das Verhalten des untersuchten Systems unter den jeweiligen Versuchsbedingungen nachvollziehbar
erfasst und anschließend mit den Ergebnissen anderer Experimente verglichen werden.

% !TeX root = ../main.tex

\begin{figure}[ht]
	\centering
	\begin{tikzpicture}
		\node[
			fill=TerminalBackground,
			draw=TerminalFrame,
			rounded corners,
			inner sep=10pt
		] {
			\begin{tikzpicture}[
				scale=0.8,
				transform shape,
				node distance=1cm,
				every node/.style={font=\small},
				startstop/.style={
					rectangle,
					rounded corners,
					draw,
					minimum width=3cm,
					minimum height=1cm,
					align=center
				},
				process/.style={
					rectangle,
					draw,
					minimum width=3cm,
					minimum height=1cm,
					align=center
				},
				decision/.style={
					diamond,
					draw,
					align=center,
					aspect=2
				},
				arrow/.style={
					->,
					>=stealth
				}
			]

			\node (start) [startstop] {Anwendung starten};
			\node (error) [process, below=of start] {Namensgebende Funktion\\ startet Experimente};
			\node (experiment) [process, below=of error] {Experiment startet Retry\\ anhand der definierten Parameter};
			\node (run) [process, below=of experiment] {SQL ausführen};
			\node (sOrf) [decision, below=of run] {success or failure};
			\node (possible) [decision, below right=1.5cm of sOrf] {Retry möglich?};
			\node (retry) [process, below right=2cm of possible] {Retry ausführen};
			\node (log) [process, below=3cm of sOrf] {Log result};
			\node (calc) [process, below=of log] {Kenngrößen berechnen};
			\node (result) [process, below=of calc] {Ergebnisse auswerten};
			\node (end) [startstop, below=of result] {Experiment beenden};

			\draw [arrow, shorten <=5pt, shorten >=5pt] (start) -- (error);
			\draw [arrow, shorten <=5pt, shorten >=5pt] (error) -- (experiment);
			\draw [arrow, shorten <=5pt, shorten >=5pt] (experiment) -- (run);
			\draw [arrow, shorten <=5pt, shorten >=5pt] (run) -- (sOrf);
			\draw [arrow, shorten <=5pt, shorten >=5pt] (sOrf) -- node[midway, above right] {failed} (possible);
			\draw [arrow, shorten <=5pt, shorten >=5pt] (sOrf) -- node[midway, above left] {succeeded} (log);
			\draw [arrow, shorten <=5pt, shorten >=5pt] (possible) -- node[below left] {possible} (retry);
			\draw [arrow, shorten <=5pt, shorten >=5pt] (possible) -- node[below right] {not possible} (log);
			\draw [arrow, shorten <=5pt, shorten >=5pt] (retry) |- (sOrf);
			\draw [arrow, shorten <=5pt, shorten >=5pt] (log) -- (calc);
			\draw [arrow, shorten <=5pt, shorten >=5pt] (calc) -- (result);
			\draw [arrow, shorten <=5pt, shorten >=5pt] (result) -- (end);

			\end{tikzpicture}};
	\end{tikzpicture}
	
	\caption{Ablauf eines Experiments}
	\label{fig:experiment_flow}
\end{figure}